\chapter{Guía de instalación}
\label{Appendix:Key1}

En este apéndice se mostrará dónde descargar la aplicación y se explicará su instalación en Linux.\\

Se ha comprobado su funcionamiento en Debian, en la primera iso descargable en:
\begin{minted}{bash}
https://cdimage.debian.org/debian-cd/current/amd64/iso-cd/
\end{minted}

Para descargar la aplicación es necesario clonar el repositorio de la herramienta, \url{https://github.com/Alberto12x/PITEA.git}. Se puede hacer de la siguiente manera:

\begin{minted}{bash}
git clone https://github.com/Alberto12x/PITEA.git
 \end{minted}
o descargando el archivo \texttt{.zip} del repositorio de \textit{GitHub}.\\

Una vez clonado, es necesario entrar en el repositorio y ejecutar el archivo \texttt{install.sh}, dándole permisos de ejecución si fuera necesario:
\begin{minted}{bash}
cd TFG-PITEA-main
chmod +x install.sh
sudo ./install.sh
\end{minted}
Una vez instalado, se llama a la herramienta por terminal escribiendo:
\begin{minted}{bash}
pitea
\end{minted}


% Caja informativa tipo aviso
\newtcolorbox{notaimportante}{
  colback=yellow!10,      % fondo amarillo claro
  colframe=yellow!50!black, % borde amarillo oscuro
  title=Nota importante,
  fonttitle=\bfseries,
  coltitle=black,
  sharp corners,
  boxrule=0.5mm,
  left=2mm,
  right=2mm,
  top=1mm,
  bottom=1mm
}

% Luego en el cuerpo del documento:
\begin{notaimportante}
Es de tener en cuenta que si se usa una \textit{shell} con el AUTO\_CD activado como zsh si se escribe en la terminal \texttt{pitea} desde la raíz del proyecto no se ejecutará la herramienta sino que se ejecutará \texttt{cd pitea}.
\end{notaimportante}
